% RESULTADO_RECUPERADO. II REV10,
% 05b_coordenadas_angulares.tex:1--125.
% Se conserva la publicación angular, sus cartas y pruebas; el funcional
% ulterior de Barbero no es una dependencia de la forma elíptica.
\section{Angular representation and channel orientation}
\label{iv:angular:section}

The section \(\hbar_{\mathrm{ret}}\) belongs to the action
line \(\mathcal L_S\); \(\eta_{\mathrm{ret}}\) is its
dimensionless coefficient in the basis specified in
\eqref{iv:action:hbar}. The angular representation then acts
on \((\alpha,\eta_{\mathrm{ret}})\). Writing
\([t]_{360}\) for the class of \(t\) in
\(\mathbb R/(360\mathbb Z)\), define
\begin{equation}
 \mathcal P_{\mathrm{ang}}(\alpha,\eta_{\mathrm{ret}})
 =\bigl([1000\alpha]_{360},
        [2(\eta_{\mathrm{ret}}+\alpha)]_{360}\bigr).
 \label{iv:angular:publicacion}
\end{equation}
Nonadic completion determines the factor \(1000\);
bidirectional closure determines the factor two. The circular
chart of \(360\) units preserves the partitions
\(12\cdot30=4\cdot90=3\cdot120\). On the positive branch, the
lifted representatives are
\begin{equation}
 A_{\deg}=1000\alpha,\qquad
 C^*_{\deg}=2(\eta_{\mathrm{ret}}+\alpha).
 \label{iv:angular:representantes}
\end{equation}
The subscript distinguishes these angular coordinates from the blocks
of the dodecaphase record.

\begin{proposition}[Shift of the base-thousand representation]
\label{iv:angular:desplazamiento}
Let \(\alpha=\sum_{n\geq1}b_n1000^{-n}\), with
\(b_n\in\{0,\ldots,999\}\), be the compatible expansion already
generated. Its lifted coordinate satisfies
\[
 A_{\deg}=b_1+\sum_{n\geq1}b_{n+1}1000^{-n}.
\]
The representation preserves the blocks from the second onward and
places the initial block in the integer part.
\end{proposition}
\begin{proof}
The series converges absolutely. Multiplying each term by
\(1000\) and separating the initial term gives the identity.
The operation acts on the constructed expansion, with its
blocks and memory; the subsequent circular representation takes
its class modulo \(360\).
\end{proof}

\subsection{Circular charts and covariance of the action}
\label{iv:angular:cartas}

\begin{definition}[Charts of the angular coordinate]
\label{iv:angular:unidades}
The degree chart expresses a complete turn by \(360\)
units; the radian chart, by \(2\pi\). The coordinate
change and its inverse are
\begin{equation}
 \theta_A=\frac{\pi}{180}A_{\deg},\qquad
 \theta_C=\frac{\pi}{180}C^*_{\deg},\qquad
 A_{\deg}=\frac{180}{\pi}\theta_A,\qquad
 C^*_{\deg}=\frac{180}{\pi}\theta_C.
 \label{iv:angular:cambio-carta}
\end{equation}
A traversal with its turn count retained uses the
lifted coordinates. The circular quotient expresses its
class modulo the corresponding turn.
\end{definition}

The record retains orientation, turns, and transport memory.
The final phase and the lift are different representations:
the latter distinguishes traversals sharing the former.
Each action expression specifies which of these coordinates it uses.

\begin{proposition}[Covariance of the action readout]
\label{iv:angular:accion}
For a section \(s\in\{\mathrm{pre},\mathrm{ret}\}\) and a
lifted angular coordinate, one has
\begin{equation}
 \mathcal S_s(\theta)
 =\hbar_s\theta_{\mathrm{rad}}
 =h_s\frac{\theta_{\deg}}{360}.
 \label{iv:angular:accion-cartas}
\end{equation}
One positive turn expresses \(h_s\).
\end{proposition}
\begin{proof}
Substituting
\(\theta_{\mathrm{rad}}=\pi\theta_{\deg}/180\) and
\(h_s=2\pi\hbar_s\) gives
\[
 \hbar_s\frac{\pi\theta_{\deg}}{180}
 =\frac{h_s\theta_{\deg}}{360}.
\]
For one turn, \(\theta_{\deg}=360\), yielding \(h_s\).
\end{proof}

The periodicity of a representation retains its own
condition: if \(U(2\pi)=-I\), operator restoration occurs
at \(4\pi\). The normalization of action per turn in
\eqref{iv:angular:accion-cartas} and the return of the represented
state thus retain their respective domains.

\subsection{Positive channels and reversal of the odd coordinate}
\label{iv:angular:canales}

In either chart, the channels are expressed through
\begin{equation}
 q_\pm=\exp[-(\theta_A\pm\theta_C)]
 =\exp\!\left[-\frac{\pi}{180}
                   (A_{\deg}\pm C^*_{\deg})\right].
 \label{iv:angular:q}
\end{equation}

\begin{proposition}[Recovery of the angular coordinates]
\label{iv:angular:inversion}
In the lifted real chart, the positive channels determine
uniquely
\begin{equation}
 A_{\deg}=-\frac{90}{\pi}\log(q_+q_-),\qquad
 C^*_{\deg}=\frac{90}{\pi}\log\frac{q_-}{q_+}.
 \label{iv:angular:recuperacion}
\end{equation}
The interchange \((q_+,q_-)\mapsto(q_-,q_+)\) preserves
\(A_{\deg}\) and reverses \(C^*_{\deg}\).
\end{proposition}
\begin{proof}
The sum and difference of the logarithms in
\eqref{iv:angular:q} are \(-\pi A_{\deg}/90\) and
\(-\pi C^*_{\deg}/90\), respectively. The real logarithm
is single-valued on the positive channels. Their product is symmetric,
and their quotient is inverted upon interchange.
\end{proof}

The antipodal involution \(r\mapsto r+6\pmod {12}\) of the
record forms six pairs. Uniform distribution of the
odd coordinate over those pairs assigns
\(C^*_{\deg}/6\) to each; their sum recovers
\(C^*_{\deg}\). Unit conversion is applied
exactly once to each argument. In particular,
\(C^*_{\deg}/A_{\deg}=\theta_C/\theta_A\) when
\(A_{\deg}\ne0\), while the preceding state retains
the lifted angles and their memory.
